\documentclass[../DoAn.tex]{subfiles}
\begin{document}
\section{Ngữ cảnh bài toán}
\label{section:2.1}
Trong lĩnh vực giao thông thông minh (ITS), bài toán nhận dạng biển số xe tự động (ALPR) đã được nghiên cứu rộng rãi và đạt được nhiều bước tiến lớn. Các giải pháp hiện hành chủ yếu tận dụng sức mạnh của các mô hình học sâu (Deep Learning) như mạng nơ-ron tích chập (CNN), tiêu biểu là họ mô hình YOLO, Faster R-CNN cho tác vụ phát hiện đối tượng, kết hợp cùng các công cụ nhận dạng ký tự quang học (OCR) như Pytesseract, PaddleOCR, CRNN hay LPRNet. Khi được tích hợp vào các nền tảng ứng dụng Web và giao tiếp qua API, các hệ thống này vận hành trơn tru và cung cấp độ chính xác rất cao trong các điều kiện môi trường lý tưởng.

Tuy nhiên, nguyên nhân chính khiến nhiều dự án và giải pháp ALPR hiện tại thất bại khi triển khai thực tiễn nằm ở sự phụ thuộc quá lớn vào các điều kiện có tính ràng buộc cao (constrained environments). Các hệ thống này thường yêu cầu camera đặt cố định ở góc chụp trực diện, phương tiện dừng đỗ hoặc di chuyển cực chậm, và ánh sáng môi trường phải đồng đều. Sự thất bại khi hoạt động trong môi trường không kiểm soát (unconstrained environments) xuất phát từ hai nguyên nhân chính. Thứ nhất, các tập dữ liệu được sử dụng để huấn luyện mô hình (như YOLO hay R-CNN) thường thiếu vắng các mẫu ảnh ở điều kiện ánh sáng yếu, lóa sáng hay góc chụp nghiêng lệch. Thứ hai, các công nghệ OCR hiện hành (điển hình như Pytesseract) vốn được thiết kế cho văn bản phẳng; chúng đòi hỏi độ tương phản đen trắng rõ rệt và ký tự không bị méo mó. Do đó, khi áp dụng lên luồng video thực tế — nơi biển số liên tục bị biến dạng phối cảnh, mờ nhòe do chuyển động và chịu nhiễu ánh sáng — khả năng trích xuất đặc trưng của mô hình giảm sút nghiêm trọng, dẫn đến việc dự đoán sai lệch hoặc bỏ sót hoàn toàn ký tự.

Từ những phân tích trên, ngữ cảnh của đồ án này được đặt ra nhằm giải quyết bài toán nhận dạng biển số đối với các phương tiện di chuyển ở tốc độ vừa phải (không quá nhanh), nhưng xuất hiện dưới các góc chụp rất đa dạng. Ngữ cảnh này mô phỏng chân thực bối cảnh thu thập dữ liệu từ camera an ninh đặt ở các góc ngã tư hoặc camera hành trình, nơi biển số xe liên tục thay đổi kích thước, vị trí và độ nghiêng theo thời gian thực. Tiêu chuẩn khắt khe nhất được đặt ra cho hệ thống là khả năng đạt đến ngưỡng nhận thức của con người: "miễn là mắt người có thể quan sát và dịch được chuỗi ký tự trên đoạn video, hệ thống cũng phải trích xuất được kết quả chính xác".

Để đạt được tiêu chuẩn "mắt người" trong một kịch bản góc chụp đa dạng như vậy, hệ thống không thể hoạt động theo cơ chế ép buộc dự đoán trên các khung hình tĩnh (single-frame) rời rạc vốn chứa đầy rủi ro nhiễu loạn. Trong thực tế, con người không bao giờ nhìn vào một khoảnh khắc tĩnh bị nhòe để đoán biển số, mà quan sát sự dịch chuyển của phương tiện qua một chuỗi thời gian để chắp vá thông tin. Do đó, bài toán đòi hỏi một luồng xử lý linh hoạt ở cấp độ video, tích hợp khả năng theo dõi quỹ đạo phương tiện và đánh giá chất lượng ảnh. Cơ chế tổng hợp thông tin dư thừa theo thời gian (temporal redundancy) cho phép hệ thống chọn lọc các khung hình có chất lượng tốt nhất trong cùng một quỹ đạo phương tiện. Nhờ đó, độ chính xác nhận dạng được cải thiện đáng kể so với các mô hình OCR xử lý từng khung hình rời rạc.

\section{Các nghiên cứu tương tự}
\label{section:2.2}
Trong những năm gần đây, bài toán nhận dạng biển số xe tự động (ALPR) đã thu hút sự chú ý mạnh mẽ từ cộng đồng nghiên cứu, đặc biệt là khi có sự ra đời của các kiến trúc học sâu. Đối với hướng tiếp cận trên các khung hình tĩnh, Pan và cộng sự đã đề xuất một hệ thống nhận dạng sử dụng mô hình YOLOv7 kết hợp với LPRNet, đạt được sự ổn định nhất định trên các tập dữ liệu biển số Trung Quốc \cite{pan2023research}. Tuy nhiên, giải pháp này nhanh chóng bộc lộ điểm yếu và suy giảm độ chính xác khi đối mặt với hiện tượng che khuất hoặc biến đổi ánh sáng phức tạp. Tương tự, Moussaoui và các cộng sự đã ứng dụng mô hình YOLOv8 kết hợp với các kỹ thuật tiền xử lý tăng cường hình ảnh để cải thiện độ rõ nét của biển số trước khi đưa vào module OCR. Dù vậy, phương pháp này vẫn chưa xử lý được phần lớn sự suy giảm tỷ lệ nhận dạng khi gặp các ký tự bị mờ nhòe do chuyển động hoặc các ký tự dính liền nhau \cite{moussaoui2024enhancing}. Ở một phạm vi hẹp hơn, Habeeb và cộng sự hay Youssef và cộng sự đã sử dụng các phiên bản YOLO cũ hơn (như YOLOv2, YOLOv3) để nhận diện biển số đặc thù của từng quốc gia (Iraq, Ai Cập). Các nghiên cứu này cho thấy tốc độ xử lý rất nhanh nhưng lại phụ thuộc quá nhiều vào kích thước đối tượng và thường chỉ hoạt động tốt trong các kịch bản bị ràng buộc chặt chẽ về khoảng cách cũng như góc chụp \cite{habeeb2021deep}.

Để giải quyết vấn đề biến dạng hình học trong các môi trường không ràng buộc, Silva và Jung đã đề xuất một luồng xử lý học sâu cho phép phát hiện và nắn chỉnh (rectification) các biển số bị méo góc trước khi đưa vào hệ thống nhận dạng. Mặc dù mang lại hiệu quả tích cực trong việc cải thiện kết quả OCR \cite{silva2018license}, bộ dữ liệu thử nghiệm của họ vẫn thiếu vắng các kịch bản thời tiết hoặc ánh sáng thực sự khắc nghiệt. Theo một hướng tiếp cận khác nhằm khắc phục giới hạn của mạng nơ-ron tích chập (CNN) truyền thống, Azadbakht và cộng sự đã áp dụng kiến trúc Vision Transformer (Vision-LPR) để nhận dạng các biển số bị mờ và có độ phân giải thấp. Mặc dù đạt được độ chính xác ấn tượng nhờ cơ chế chú ý toàn cục (global attention), phương pháp này lại đòi hỏi chi phí tính toán lớn và có tốc độ suy luận chậm, gây rào cản lớn cho việc triển khai trên các hệ thống giám sát yêu cầu thời gian thực \cite{azadbakht2022multipath}.

Đáng chú ý nhất, đối với kịch bản nhận dạng phương tiện di chuyển trên luồng dữ liệu video, một số nghiên cứu đã bắt đầu nhận ra rủi ro của việc dựa vào khung hình đơn lẻ và chuyển hướng sang khai thác đặc trưng chuỗi thời gian. Nhóm tác giả Laroca, Severo và cộng sự \cite{laroca2018robust}
 đã giới thiệu hệ thống ALPR thời gian thực sử dụng mạng Fast-YOLO, đồng thời ra mắt bộ dữ liệu UFPR-ALPR mô phỏng kịch bản cả camera và phương tiện cùng di chuyển. Cải tiến đáng chú ý trong nghiên cứu của họ là việc bước đầu áp dụng thông tin dư thừa theo thời gian (temporal redundancy). Bằng cách nhận dạng độc lập từng khung hình của cùng một phương tiện và áp dụng cơ chế bầu chọn theo số đông (majority voting), tỷ lệ nhận dạng của họ đã tăng vọt từ 85.45\% lên 93.53\% trên tập dữ liệu SSIG. Mặc dù cơ chế này đã chứng minh được sức mạnh của việc kết hợp đa khung hình, các phương pháp bầu chọn nêu trên vẫn mang tính chất thống kê cơ bản. Chúng thiếu vắng một thuật toán theo vết đa đối tượng (multi-object tracking) độ chính xác cao và chưa có bộ định tuyến để đánh giá chất lượng (quality routing) nhằm chủ động loại bỏ các khung hình nhiễu trước khi tiến hành đối sánh. Đây chính là khoảng trống công nghệ lớn nhất mà đồ án này tập trung khai thác và hoàn thiện.
\section{Các mô hình phát hiện đối tượng}
\label{section:2.3}
\subsection{YOLOv5}
\label{subsection:2.3.1}
YOLOv5 là một họ mô hình phát hiện đối tượng một giai đoạn, trực tiếp dự đoán đồng thời hộp bao và nhãn lớp từ bản đồ đặc trưng mà không cần giai đoạn sinh vùng đề xuất như các mô hình hai giai đoạn (Faster R-CNN, Mask R-CNN). Nhờ đó, độ trễ suy luận được giảm thiểu đáng kể, phù hợp với các hệ thống xử lý video thời gian thực.

Trong đồ án này, YOLOv5 được kế thừa từ Che et al.\ \cite{che2023ctm} và sử dụng nguyên vẹn không qua fine-tune, đảm nhiệm duy nhất tầng phát hiện phương tiện. Đầu ra là các hộp bao quanh xe máy, ô tô, xe tải trong khung hình, phục vụ làm vùng quan tâm (RoI) cho tầng phát hiện biển số phía sau.

Mô hình phù hợp với vai trò này vì ba lý do chính. Thứ nhất, kiến trúc một giai đoạn mang lại tốc độ cao, đáp ứng yêu cầu xử lý liên tục trên luồng video. Thứ hai, cơ chế neck PANet tổng hợp đặc trưng đa tỉ lệ, cho phép phát hiện phương tiện ở nhiều kích thước khác nhau — từ xe máy nhỏ ở xa đến xe tải lớn ở gần camera — trong cùng một khung hình. Thứ ba, định vị phương tiện trước giúp pipeline hoạt động phân cấp: khoanh vùng tìm kiếm, loại bỏ nhiễu nền và cung cấp vùng ảnh sạch hơn cho tầng phát hiện biển số.

\begin{figure}[htbp]
    \centering
    \includegraphics[width=1\textwidth]{Hinhve/yolov5_architecture.png}
    \caption{Kiến trúc tổng quát của mô hình YOLOv5 \cite{nagvekar2023yolov5}}
    \label{fig:yolov5_architecture}
\end{figure}

\subsection{YOLOv8 OBB}
\label{subsection:2.3.2}
Bên cạnh kiến trúc phát hiện đối tượng tiêu chuẩn, YOLOv8 cung cấp thêm biến thể YOLOv8 OBB \cite{ultralytics_yolov8} được thiết kế đặc thù để nhận diện các vật thể có độ nghiêng và góc xoay phức tạp. Điểm khác biệt chính trong kiến trúc của YOLOv8 OBB so với các mô hình truyền thống nằm ở phần đầu dự đoán phân tách không neo (anchor-free split Ultralytics head). Thay vì bị giới hạn bởi việc dự đoán các hộp giới hạn song song với trục tọa độ của bức ảnh, mạng nơ-ron được bổ sung thêm một tham số góc xoay ở lớp đầu ra, chuyển đổi không gian tính toán sang định dạng xywhr. Trong cấu trúc biểu diễn này, thành phần x và y đại diện cho tọa độ tâm, w và h là chiều rộng và chiều cao, còn thành phần r đóng vai trò lưu trữ góc xoay của đối tượng so với trục tham chiếu. Thông qua cơ chế lan truyền thẳng, mạng sẽ xuất ra véc-tơ năm chiều này, sau đó hệ thống suy luận mới thực hiện quá trình ánh xạ toán học để chuyển đổi từ định dạng xywhr về tọa độ không gian chuẩn hóa của bốn điểm góc.

Sự thay đổi về biểu diễn tham số hình học đòi hỏi hệ thống hàm mất mát (loss function) của YOLOv8 OBB cũng phải được thiết kế lại. Quá trình tính toán lỗi định vị không chỉ dựa trên mức độ chồng lấp của hai hình chữ nhật thông thường, mà mọi sai số đều được xử lý trực tiếp trên không gian xywhr. Việc kết hợp tham số góc xoay vào hàm mất mát IoU (Intersection over Union) giúp mạng nơ-ron bị phạt nặng hơn khi dự đoán sai hướng của đối tượng, buộc mô hình phải học cách tinh chỉnh góc xoay sao cho đường viền của hộp giới hạn ôm sát nhất vào các cạnh thực tế của vật thể. Đặc điểm kiến trúc không neo (anchor-free) kết hợp cùng hàm mất mát tối ưu góc xoay này giúp mô hình triệt tiêu hoàn toàn sự phụ thuộc vào các hộp neo cố định, mang lại độ linh hoạt cao khi đối tượng bị xoay ở các biên độ bất kỳ.

Đối với bài toán nhận dạng biển số xe trên luồng video không ràng buộc, phương tiện giao thông thường xuyên di chuyển với quỹ đạo đa dạng, khiến hình ảnh biển số thu được hiếm khi nằm ngang chuẩn. Việc ứng dụng cấu trúc đầu ra xywhr của YOLOv8 OBB cho phép hệ thống xác định chuẩn xác tọa độ bốn điểm góc của biển số bất chấp các góc nghiêng lệch. Dữ liệu tọa độ không gian này là cơ sở hình học then chốt để áp dụng phép toán biến đổi phối cảnh (perspective transform). Quá trình này sẽ kéo phẳng vùng biển số bị méo về dạng hình chữ nhật chuẩn, cắt bỏ tối đa các điểm ảnh thuộc về cảnh nền vốn thường xuất hiện ở bốn góc nếu dùng hộp giới hạn thẳng đứng. Nhờ đặc tính kiến trúc vượt trội này, hệ thống cung cấp các vùng ảnh đầu vào được chuẩn hóa hình học, giúp nâng cao độ ổn định và chính xác cho mô hình nhận dạng chuỗi ký tự (OCR) ở giai đoạn tiếp theo.

\begin{figure}[htbp]
    \centering
    \includegraphics[width=0.9\textwidth]{Hinhve/hbb_vs_obb_comparison_final.jpg}
    \caption{So sánh trực quan giữa khả năng định vị của hộp bao ngang (HBB) và hộp bao có hướng (OBB) trên biển số xe}
    \label{fig:hbb_vs_obb}
\end{figure}

\section{Thuật toán theo vết đối tượng BoT-SORT}
\label{section:2.4}
BoT-SORT (Robust Associations Multi-Pedestrian Tracking) là thuật toán bám vết đa đối tượng trực tuyến được Aharon, Orfaig và Bobrovsky đề xuất trong bối cảnh các tracker họ SORT vẫn còn hạn chế về mô hình chuyển động, bù chuyển động camera và cách kết hợp đặc trưng ngoại hình \cite{aharon2022bot}. Thuật toán đi theo mô hình \textit{tracking-by-detection}: tại mỗi khung hình, bộ phát hiện đối tượng sinh ra các hộp bao mới; bộ theo dõi dự đoán trạng thái của các track đã tồn tại; sau đó bài toán gán giữa track và detection được giải như một bài toán tối ưu ghép cặp. BoT-SORT không thay thế detector, mà tập trung cải thiện ba thành phần của tracker: bộ lọc Kalman, bù chuyển động camera và hàm chi phí kết hợp IoU--ReID.

Với mô hình chuyển động, BoT-SORT vẫn dùng bộ lọc Kalman rời rạc dưới giả thiết vận tốc không đổi, nhưng thay đổi cách biểu diễn trạng thái hộp bao. Trong SORT, trạng thái dùng diện tích hộp và tỷ lệ khung hình; trong DeepSORT, trạng thái dùng tỷ lệ khung hình \(a\) và chiều cao \(h\). Theo các tác giả, cách biểu diễn này có thể làm ước lượng chiều rộng thiếu chính xác. Vì vậy BoT-SORT định nghĩa trực tiếp trạng thái theo tâm, chiều rộng, chiều cao và đạo hàm theo thời gian:
\[
\mathbf{x}_k =
[x_c(k), y_c(k), w(k), h(k),
\dot{x}_c(k), \dot{y}_c(k), \dot{w}(k), \dot{h}(k)]^\top .
\]
Vector đo từ detector chỉ gồm bốn đại lượng quan sát được:
\[
\mathbf{z}_k =
[z_{x_c}(k), z_{y_c}(k), z_w(k), z_h(k)]^\top .
\]
Với ma trận chuyển trạng thái \(\mathbf{F}\), ma trận quan sát \(\mathbf{H}\), hiệp phương sai trạng thái \(\mathbf{P}\), nhiễu quá trình \(\mathbf{Q}_k\) và nhiễu đo \(\mathbf{R}_k\), bước dự đoán có dạng:
\[
\hat{\mathbf{x}}_{k|k-1}=\mathbf{F}\hat{\mathbf{x}}_{k-1|k-1},
\qquad
\mathbf{P}_{k|k-1}=
\mathbf{F}\mathbf{P}_{k-1|k-1}\mathbf{F}^{\top}+\mathbf{Q}_k.
\]
Khi có detection được gán vào track, Kalman gain và trạng thái hậu nghiệm được cập nhật bởi:
\[
\mathbf{K}_k =
\mathbf{P}_{k|k-1}\mathbf{H}^{\top}
(\mathbf{H}\mathbf{P}_{k|k-1}\mathbf{H}^{\top}+\mathbf{R}_k)^{-1},
\]
\[
\hat{\mathbf{x}}_{k|k}=
\hat{\mathbf{x}}_{k|k-1}
+\mathbf{K}_k(\mathbf{z}_k-\mathbf{H}\hat{\mathbf{x}}_{k|k-1}).
\]
Trong paper, \(\mathbf{Q}_k\) và \(\mathbf{R}_k\) được đặt phụ thuộc vào kích thước hộp bao, ví dụ các nhiễu vị trí và kích thước tỷ lệ với \(\hat{w}\) và \(\hat{h}\). Cách đặt này phản ánh trực giác rằng hộp lớn thường có sai số tuyệt đối lớn hơn hộp nhỏ.

Thành phần thứ hai của BoT-SORT là bù chuyển động camera (Camera Motion Compensation). Trong các tracker dựa trên IoU, nếu camera dịch chuyển, hộp dự đoán của track có thể lệch khỏi detection thật dù đối tượng không đổi danh tính. BoT-SORT ước lượng chuyển động nền giữa hai khung hình bằng đăng ký ảnh và biểu diễn nó bằng ma trận affine:
\[
\mathbf{A}_{k-1}^{k} =
[\mathbf{M}\mid \mathbf{T}] =
\begin{bmatrix}
a_{11} & a_{12} & a_{13}\\
a_{21} & a_{22} & a_{23}
\end{bmatrix}.
\]
Từ \(\mathbf{M}\) và \(\mathbf{T}\), thuật toán xây dựng ma trận mở rộng \(\tilde{\mathbf{M}}\) và vector \(\tilde{\mathbf{T}}\) để hiệu chỉnh dự đoán của Kalman:
\[
\hat{\mathbf{x}}'_{k|k-1} =
\tilde{\mathbf{M}}\hat{\mathbf{x}}_{k|k-1}+\tilde{\mathbf{T}},
\qquad
\mathbf{P}'_{k|k-1} =
\tilde{\mathbf{M}}\mathbf{P}_{k|k-1}\tilde{\mathbf{M}}^{\top}.
\]
Sau bước này, trạng thái đã được đưa về hệ tọa độ của khung hình hiện tại trước khi ghép nối với detection.

Thành phần thứ ba là chiến lược kết hợp IoU và ReID. Với track \(i\) và detection \(j\), khoảng cách hình học được tính bằng:
\[
d^{iou}_{ij}=1-\operatorname{IoU}(\hat{b}_i,b_j),
\]
trong đó \(\hat{b}_i\) là hộp dự đoán của track và \(b_j\) là hộp detection. Khoảng cách ngoại hình là cosine distance giữa đặc trưng trung bình của track \(e_i\) và đặc trưng detection \(f_j\):
\[
d^{cos}_{ij}=1-e_i^\top f_j.
\]
Đặc trưng track được cập nhật bằng trung bình trượt mũ:
\[
e_i^k=\alpha e_i^{k-1}+(1-\alpha)f_i^k.
\]
Thay vì dùng tổng có trọng số \(C=\lambda A_a+(1-\lambda)A_m\), BoT-SORT dùng cơ chế lọc theo ngưỡng rồi lấy chi phí nhỏ hơn:
\[
\hat{d}^{cos}_{ij}=
\begin{cases}
0.5\,d^{cos}_{ij}, & d^{cos}_{ij}<\theta_{emb}\ \land\ d^{iou}_{ij}<\theta_{iou},\\
1, & \text{ngược lại},
\end{cases}
\qquad
C_{ij}=\min\{d^{iou}_{ij},\hat{d}^{cos}_{ij}\}.
\]
Ma trận chi phí \(C\) được đưa vào thuật toán Hungarian để tìm phép gán tối ưu. Như vậy, BoT-SORT giữ được ưu điểm của IoU trong định vị ngắn hạn, đồng thời tận dụng ReID để duy trì định danh khi các đối tượng gần nhau, bị che khuất hoặc có chuyển động phức tạp.

\section{Các mô hình nhận dạng kí tự quang học (OCR)}
\label{section:2.5}

\subsection{Các phương pháp nhận dạng phân đoạn ký tự}
\label{subsection:2.5.1}
Các phương pháp nhận dạng biển số xe dựa trên phân đoạn ký tự (segmentation-based) hoạt động theo nguyên lý tách rời bước định vị ký tự ra khỏi bước nhận dạng. Sau khi ảnh biển số được trích xuất, hệ thống thực hiện phân đoạn để xác định vùng ảnh chứa từng ký tự riêng lẻ, sau đó đưa từng ký tự đã phân đoạn vào bộ phân loại hoặc bộ nhận dạng độc lập. Do đó, chất lượng của toàn bộ hệ thống phụ thuộc trực tiếp vào độ chính xác của giai đoạn phân đoạn. Mục này tập trung vào hai nhóm phương pháp phân đoạn ký tự có tính đại diện trong giai đoạn phát triển từ học máy truyền thống sang học sâu, bao gồm: phương pháp dựa trên trích xuất đặc trưng thủ công kết hợp phân loại thống kê, và phương pháp phát hiện đối tượng bằng mạng nơ-ron tích chập sâu.

\paragraph{Phương pháp học máy truyền thống}

Trước khi các kiến trúc mạng nơ-ron tích chập sâu được ứng dụng rộng rãi, nhóm phương pháp học máy truyền thống là hướng tiếp cận chủ đạo trong bài toán nhận dạng ký tự biển số. Trong các phương pháp này, mỗi ký tự sau khi được phân đoạn ra khỏi ảnh biển số sẽ được đưa qua một pipeline gồm ba bước tuần tự: trích xuất đặc trưng thủ công, giảm chiều dữ liệu, và phân loại \cite{liu2012research}.

Về trích xuất đặc trưng, Histogram of Oriented Gradients (HOG) là descriptor thủ công phổ biến. Phương pháp này hoạt động bằng cách chia ảnh ký tự thành các ô (cell) cố định, tính histogram gradient theo nhiều hướng trong từng ô, sau đó chuẩn hóa theo khối (block) để tăng tính bất biến với thay đổi độ sáng cục bộ. Khan et al. đề xuất kiến trúc kết hợp kênh độ sáng L từ không gian màu CIE-Lab, nhị phân hóa thích nghi, sau đó hợp nhất đặc trưng HOG với các đặc trưng hình học (geometric features), và áp dụng phương pháp lựa chọn đặc trưng dựa trên entropy trước khi đưa vào bộ phân loại SVM. Để giảm chiều không gian đặc trưng và giảm số chiều và giữ lại các thành phần chính của dữ liệu, nhiều nghiên cứu áp dụng thêm Phân tích thành phần chính (PCA) trước bước phân loại. Về phía phân loại, Support Vector Machine (SVM) được huấn luyện theo phương thức one-vs-all để nhận dạng 36 lớp ký tự (10 chữ số và 26 chữ cái). Một hướng khác là sử dụng Hidden Markov Model (HMM). Trong mô hình HMM, mỗi ảnh ký tự sau khi được biểu diễn thành vector đặc trưng trực quan, chẳng hạn HOG, được xem là một quan sát. Từ chuỗi quan sát này, HMM suy luận chuỗi trạng thái ẩn tương ứng với các ký tự hoặc các vị trí cấu trúc trong biển số.

Nhược điểm căn bản của nhóm phương pháp này nằm ở việc phụ thuộc hoàn toàn vào thiết kế đặc trưng thủ công, vốn thiếu khả năng học các biểu diễn phức tạp và thích nghi tự động với dữ liệu trong môi trường thực tế. Quan trọng hơn, toàn bộ hiệu năng của hệ thống phụ thuộc nghiêm ngặt vào chất lượng của bước phân đoạn ký tự đứng trước: nếu bước phân đoạn tạo ra ký tự bị cắt thiếu, dính liền hay lệch vị trí, bộ phân loại SVM hay HMM hoàn toàn không có cơ chế bù đắp sai sót này.

\paragraph{Phương pháp phát hiện ký tự bằng mô hình học sâu}
Laroca et al. \cite{laroca2018robust} là một trong những người tiên phong theo hướng này, ông cùng cộng sự đề xuất hệ thống ALPR hai giai đoạn hoàn toàn dựa trên kiến trúc YOLO. Cụ thể, các tác giả sử dụng Fast-YOLO để phát hiện phương tiện và định vị biển số, sau đó áp dụng một kiến trúc lấy cảm hứng từ Fast-YOLO thứ hai — được các tác giả đặt tên là CR-NET, bao gồm 11 lớp đầu của YOLO cùng bốn lớp tích chập bổ sung để tăng phi tuyến tính — chuyên biệt cho bài toán phân đoạn và nhận dạng ký tự đồng thời. Mô hình CR-NET xử lý ảnh biển số đã được cắt ra, dự đoán bounding box và nhãn lớp ký tự trong một lần truyền thuận, và thứ tự chuỗi được suy luận theo tọa độ hoành độ của tâm bounding box. Hệ thống đạt tỷ lệ nhận dạng 93.53\% và tốc độ 47 FPS trên tập dữ liệu SSIG. Một điểm đáng chú ý là các tác giả sử dụng kỹ thuật bù dư thừa thời gian (temporal redundancy) bằng cách tổng hợp dự đoán từ nhiều khung hình liên tiếp tương ứng với cùng một phương tiện theo phương thức bỏ phiếu đa số, nhờ đó cải thiện tỷ lệ nhận dạng từ 85.45\% lên 93.53\%.

Selmi et al. \cite{selmi2020delpdar} đề xuất DELP-DAR, một hệ thống dựa trên Mask R-CNN để phát hiện biển số, phát hiện ký tự và nhận dạng ký tự. Hệ thống sử dụng backbone tương tự GoogLeNet được rút gọn nhằm giảm chi phí tính toán, kết hợp các luật hậu xử lý để lọc vùng ký tự không hợp lệ. Về kết quả, phương pháp đạt precision và recall lần lượt 98.9\% và 98.6\% trong phát hiện biển số trên Caltech; trên PKU, tỷ lệ phát hiện trung bình đạt 99.4\%. Ở bước nhận dạng ký tự, DELP-DAR đạt 97.2\% trên Caltech, từ 96.3\% đến 97.8\% trên các tập con của AOLP, và 96.9\%--97.5\% trên hai tập dữ liệu biển số Tunisia.


Mặc dù các phương pháp phát hiện ký tự dựa trên học sâu cho thấy ưu thế vượt trội so với các kỹ thuật truyền thống về độ chính xác trong điều kiện môi trường phức tạp, chúng vẫn tồn tại một hạn chế chung đáng chú ý: để huấn luyện hiệu quả, các mô hình này đòi hỏi lượng dữ liệu lớn với nhãn bounding box được gán thủ công ở cấp độ từng ký tự riêng lẻ — một công việc tốn kém và mất nhiều thời gian hơn so với gán nhãn ở cấp độ chuỗi biển số. Ngoài ra, thứ tự của các ký tự trong chuỗi kết quả phải được suy luận gián tiếp từ tọa độ bounding box chứ không được mô hình hóa trực tiếp, điều này dễ gây lỗi với các biển số có bố cục phi tuyến tính hoặc nhiều hàng.

\subsection{Các phương pháp nhận dạng không phân đoạn ký tự}
\label{subsection:2.5.2}
Nhận dạng chuỗi ký tự không cần phân đoạn (segmentation-free) đã trở thành tiêu chuẩn mới trong các hệ thống nhận dạng văn bản trong tự nhiên cũng như bài toán ALPR hiện đại. Phương pháp này khắc phục hiệu quả nhược điểm của các kỹ thuật truyền thống vốn dễ thất bại khi các ký tự dính liền nhau, bị che khuất hoặc mờ nhòe do chuyển động. Dựa trên cơ chế giải mã và mô hình hóa chuỗi, các phương pháp này được chia thành hai nhóm kiến trúc chính, cùng với sự ra đời của các nền tảng đóng gói phục vụ triển khai thực tiễn.

\paragraph{Nhóm mô hình dựa trên hàm mất mát CTC (CTC-based Models)}
\mbox{}\\
Connectionist Temporal Classification (CTC) là một kỹ thuật quan trọng cho phép mạng nơ-ron ánh xạ trực tiếp từ chuỗi đặc trưng hình ảnh sang chuỗi ký tự đích mà không yêu cầu dữ liệu phải được phân mảnh sẵn ở cấp độ từng ký tự. Kiến trúc tiêu biểu cho nhóm này là CRNN \cite{shi2017crnn}, kết hợp CNN để trích xuất đặc trưng ảnh, các lớp BiLSTM để mô hình hóa phụ thuộc theo chiều ngang của chuỗi, và CTC cho bước phiên mã. Trong quá trình giải mã, CTC loại bỏ nhãn rỗng và gộp các nhãn lặp liên tiếp để tạo ra chuỗi ký tự cuối cùng \cite{shi2017crnn}.

Tuy nhiên, trong các bài toán yêu cầu tốc độ thời gian thực như nhận dạng biển số, việc sử dụng các lớp LSTM thường tạo ra điểm nghẽn về mặt tính toán. Để giải quyết vấn đề này, các kiến trúc như LPRNet \cite{zherzdev2018lprnet} đã được đề xuất nhằm loại bỏ hoàn toàn mạng RNN, chỉ sử dụng mạng CNN nhẹ kết hợp trực tiếp với hàm mất mát CTC. Ưu điểm chính của họ mô hình dựa trên CTC là khả năng xuất ra toàn bộ chuỗi ký tự chỉ thông qua một lần lan truyền tiến (single forward pass) duy nhất, loại bỏ vòng lặp giải mã phức tạp. Gần đây, việc kết hợp CTC với các mô hình thị giác tiên tiến như SVTRv2 \cite{gu2022svtr} cũng cho thấy tốc độ và độ chính xác vượt trội so với các kiến trúc mã hóa - giải mã (Encoder-Decoder) phức tạp. Đây là hướng tiếp cận đặc biệt phù hợp cho các luồng xử lý video đòi hỏi độ trễ cực thấp.

\paragraph{Nhóm mô hình dựa trên cơ chế Chú ý và Transformer (Attention/Transformer-based Models)}
\mbox{}\\
Để giải quyết bài toán nhận dạng văn bản bị méo mó, biến dạng phối cảnh hoặc che khuất nghiêm trọng, các kiến trúc dựa trên cơ chế tự chú ý (Self-Attention) và Transformer đã được ứng dụng rộng rãi nhờ khả năng nắm bắt ngữ cảnh toàn cục mạnh mẽ. Tiêu biểu, mô hình TrOCR \cite{li2021trocr} loại bỏ hoàn toàn mạng CNN truyền thống, thay vào đó tận dụng sức mạnh của các mạng Transformer thị giác (ViT, DeiT) và mạng ngôn ngữ (RoBERTa) đã được huấn luyện trước ở quy mô lớn để sinh văn bản trực tiếp theo từng mảnh từ vựng (wordpiece).

Tuy nhiên, các mô hình sinh tự hồi quy (Autoregressive - AR) thông thường thường bị giới hạn bởi tính đơn hướng (từ trái sang phải) và tốc độ giải mã chậm. Để khắc phục, kiến trúc PARSeq \cite{bautista2022scene} đã áp dụng mô hình ngôn ngữ hoán vị (Permutation Language Modeling), cho phép mạng nơ-ron học được các xác suất có điều kiện dựa trên các tập con ngữ cảnh tùy ý. Thiết kế này hợp nhất khả năng giải mã không tự hồi quy (NAR) song song và tinh chỉnh lặp lại (iterative refinement), mang lại độ ổn định cao đối với các văn bản bị xoay nghiêng hoặc mờ nhòe trong tự nhiên. Dù đem lại năng lực mô hình hóa ngôn ngữ vượt trội, hạn chế lớn của các mô hình họ Transformer là dung lượng tham số lớn, chi phí tính toán cao và có chi phí suy luận cao hơn so với các mô hình CNN nhẹ. Việc áp dụng chúng vào bài toán biển số xe — vốn có cấu trúc ngữ pháp rất ngắn và đơn giản — đòi hỏi sự đánh đổi khắt khe về tài nguyên hệ thống.

\paragraph{Các framework OCR thực tiễn (Practical OCR Frameworks) }
\mbox{}\\
Song song với các nghiên cứu lý thuyết, việc triển khai công nghệ OCR vào các sản phẩm thực tế đã được thúc đẩy mạnh mẽ bởi các nền tảng mã nguồn mở tối ưu hóa cao. Nổi bật nhất là EasyOCR \cite{jaided2020easyocr}, một framework cung cấp giải pháp toàn trình sử dụng thuật toán CRAFT cho việc phát hiện và mạng CRNN (ResNet + LSTM + CTC) để nhận dạng. Nền tảng này hỗ trợ hơn 80 ngôn ngữ và cho phép dễ dàng tích hợp, thử nghiệm nhanh chóng trên các phần cứng phổ thông.
Ở mức độ tối ưu công nghiệp sâu sắc hơn, hệ sinh thái PaddleOCR \cite{paddleocr2020} cung cấp các pipeline (luồng xử lý) như PP-OCRv4 với hai phiên bản Server và Mobile. Bằng việc chắt lọc các kỹ thuật học sâu tiên tiến nhất, các mô hình này đạt được sự cân bằng hiệu quả giữa dung lượng nhỏ gọn, tốc độ suy luận rất nhanh (chỉ khoảng vài mili-giây trên GPU) và độ chính xác cao. Ngược lại, các công cụ truyền thống như Pytesseract OCR, dù đã được nâng cấp lõi nhận dạng LSTM, vẫn bộc lộ nhiều rào cản khi ứng dụng vào hệ thống giao thông. Do bản chất được thiết kế cho văn bản phẳng (document OCR), chúng đòi hỏi điều kiện ánh sáng chuẩn, độ tương phản đen trắng rõ rệt và thường thường thất bại trước các biển số bị biến dạng phối cảnh, nhiễu hạt hoặc thiếu sáng trong các kịch bản môi trường không kiểm soát. Sự đối lập này một lần nữa khẳng định sự cần thiết của các luồng xử lý nhận dạng chuyên biệt tích hợp sâu công nghệ học sâu thay vì phụ thuộc vào các công cụ OCR đại trà.

%%%%%%%%%%%%%%%%%%%%%%%%%%%%%%%%%%%

\end{document}
